% GENERATED FILE. Edit canonical lessons, then run npm run generate:books.
% canonical-source-hash: fnv1a64:f62d4a3045bfff05
% canonical-lessons: SA-C165-dhulih, SA-C165-dhumah, SA-C165-jvala, SA-C165-angarah, SA-C165-lauham

\chapter{Dust, Smoke, Flame, Live coal, Iron}
\label{ch:sa-dust-smoke-flame-live-coal-iron}
\hlchaptermodality{165}

\begin{chapteropening}
\textbf{By the end of this chapter:} \emph{I can say dust, smoke, a flame, a live coal and iron.}

Complete the last lesson of chapter 165: I can say dust, smoke, a flame, a live coal and iron.
\end{chapteropening}

\section[\texorpdfstring{\sk{धूलिः} (dhūliḥ)}{धूलिः (dhūliḥ)}]{\sk{धूलिः} (dhūliḥ) — dust}

\label{lesson:SA-C165-dhulih}



\begin{quote}
Before the new one: say the Sanskrit for a playground, then the Sanskrit for an airport.
\end{quote}

\subsection*{You'll want to know: \sk{धूलिः}}
\textbf{\sk{धूलिः}} — \emph{dhūliḥ} — \textquotedblleft{}dust\textquotedblright{}.

\begin{grammarlens}[title={What you've built}]
One more word for this chapter.
\end{grammarlens}

\subsection*{Your turn}
\begin{itemize}
\raggedright
  \item \emph{Say it:} \emph{dhūliḥ}
  \item \emph{Say it:} \emph{dhūliḥ}, once more
  \item \emph{Say it:} say \emph{dhūliḥ}
  \item \emph{Recall:} say the Sanskrit for a playground, then the Sanskrit for an airport, then say \emph{dhūliḥ} again
\end{itemize}

\begin{tcolorbox}[breakable,colback=teal!4,colframe=teal!35!black,title={Before you move on}]
What does \sk{धूलिः} mean? (\textquotedblleft{}Dust\textquotedblright{}.) Say it once more.
\end{tcolorbox}
\section[\texorpdfstring{\sk{धूमः} (dhūmaḥ)}{धूमः (dhūmaḥ)}]{\sk{धूमः} (dhūmaḥ) — smoke}

\label{lesson:SA-C165-dhumah}



\begin{quote}
Before the new one: say the Sanskrit for an airport, then the Sanskrit for dust.
\end{quote}

\subsection*{You'll want to know: \sk{धूमः}}
\textbf{\sk{धूमः}} — \emph{dhūmaḥ} — \textquotedblleft{}smoke\textquotedblright{}.

\begin{grammarlens}[title={What you've built}]
One more word for this chapter.
\end{grammarlens}

\subsection*{Your turn}
\begin{itemize}
\raggedright
  \item \emph{Say it:} \emph{dhūmaḥ}
  \item \emph{Say it:} \emph{dhūmaḥ}, once more
  \item \emph{Say it:} say \emph{dhūmaḥ}
  \item \emph{Recall:} say the Sanskrit for an airport, then the Sanskrit for dust, then say \emph{dhūmaḥ} again
\end{itemize}

\begin{tcolorbox}[breakable,colback=teal!4,colframe=teal!35!black,title={Before you move on}]
What does \sk{धूमः} mean? (\textquotedblleft{}Smoke\textquotedblright{}.) Say it once more.
\end{tcolorbox}
\section[\texorpdfstring{\sk{ज्वाला} (jvālā)}{ज्वाला (jvālā)}]{\sk{ज्वाला} (jvālā) — a flame}

\label{lesson:SA-C165-jvala}



\begin{quote}
Before the new one: say the Sanskrit for dust, then the Sanskrit for smoke.
\end{quote}

\subsection*{You'll want to know: \sk{ज्वाला}}
\textbf{\sk{ज्वाला}} — \emph{jvālā} — \textquotedblleft{}a flame\textquotedblright{}.

\begin{grammarlens}[title={What you've built}]
One more word for this chapter.
\end{grammarlens}

\subsection*{Your turn}
\begin{itemize}
\raggedright
  \item \emph{Say it:} \emph{jvālā}
  \item \emph{Say it:} \emph{jvālā}, once more
  \item \emph{Say it:} say \emph{jvālā}
  \item \emph{Recall:} say the Sanskrit for dust, then the Sanskrit for smoke, then say \emph{jvālā} again
\end{itemize}

\begin{tcolorbox}[breakable,colback=teal!4,colframe=teal!35!black,title={Before you move on}]
What does \sk{ज्वाला} mean? (\textquotedblleft{}A flame\textquotedblright{}.) Say it once more.
\end{tcolorbox}
\section[\texorpdfstring{\sk{अंगारः} (aṅgāraḥ)}{अंगारः (aṅgāraḥ)}]{\sk{अंगारः} (aṅgāraḥ) — a live coal}

\label{lesson:SA-C165-angarah}



\begin{quote}
Before the new one: say the Sanskrit for smoke, then the Sanskrit for a flame.
\end{quote}

\subsection*{You'll want to know: \sk{अंगारः}}
\textbf{\sk{अंगारः}} — \emph{aṅgāraḥ} — \textquotedblleft{}a live coal\textquotedblright{}.

\begin{grammarlens}[title={What you've built}]
One more word for this chapter.
\end{grammarlens}

\subsection*{Your turn}
\begin{itemize}
\raggedright
  \item \emph{Say it:} \emph{aṅgāraḥ}
  \item \emph{Say it:} \emph{aṅgāraḥ}, once more
  \item \emph{Say it:} say \emph{aṅgāraḥ}
  \item \emph{Recall:} say the Sanskrit for smoke, then the Sanskrit for a flame, then say \emph{aṅgāraḥ} again
\end{itemize}

\begin{tcolorbox}[breakable,colback=teal!4,colframe=teal!35!black,title={Before you move on}]
What does \sk{अंगारः} mean? (\textquotedblleft{}A live coal\textquotedblright{}.) Say it once more.
\end{tcolorbox}
\section[\texorpdfstring{\sk{लौहम्} (lauham)}{लौहम् (lauham)}]{\sk{लौहम्} (lauham) — iron}

\label{lesson:SA-C165-lauham}



\begin{quote}
Before the new one: say the Sanskrit for a flame, then the Sanskrit for a live coal.
\end{quote}

\subsection*{You'll want to know: \sk{लौहम्}}
\textbf{\sk{लौहम्}} — \emph{lauham} — \textquotedblleft{}iron\textquotedblright{}.

\begin{grammarlens}[title={What you've built}]
One more word for this chapter.
\end{grammarlens}

\subsection*{Your turn}
\begin{itemize}
\raggedright
  \item \emph{Say it:} \emph{lauham}
  \item \emph{Say it:} \emph{lauham}, once more
  \item \emph{Say it:} say \emph{lauham}
  \item \emph{Recall:} say the Sanskrit for a flame, then the Sanskrit for a live coal, then say \emph{lauham} again
\end{itemize}

\begin{tcolorbox}[breakable,colback=teal!4,colframe=teal!35!black,title={Before you move on}]
What does \sk{लौहम्} mean? (\textquotedblleft{}Iron\textquotedblright{}.) Say it once more.
\end{tcolorbox}
